%% University Curriculum Flow and Capacity Planning
%% IEEE conference format. Compile on Overleaf (pdfLaTeX) or locally with IEEEtran installed.
%% >>> REPLACE the four author blocks below with the real names and roll numbers. <<<

\documentclass[conference]{IEEEtran}
\IEEEoverridecommandlockouts

\usepackage{cite}
\usepackage{amsmath,amssymb,amsfonts}
\usepackage{algorithmic}
\usepackage{graphicx}
\usepackage{textcomp}
\usepackage{xcolor}
\usepackage{booktabs}
\usepackage{url}
\def\BibTeX{{\rm B\kern-.05em{\sc i\kern-.025em b}\kern-.08em T\kern-.1667em\lower.7ex\hbox{E}\kern-.125emX}}

\begin{document}

\title{University Curriculum Flow and Capacity Planning\\
Using Integer and Binary Programming}

\author{
\IEEEauthorblockN{Author One}
\IEEEauthorblockA{\textit{Department} \\
\textit{University}\\
City, Country \\
ROLL.NUMBER.ONE}
\and
\IEEEauthorblockN{Author Two}
\IEEEauthorblockA{\textit{Department} \\
\textit{University}\\
City, Country \\
ROLL.NUMBER.TWO}
\and
\IEEEauthorblockN{Author Three}
\IEEEauthorblockA{\textit{Department} \\
\textit{University}\\
City, Country \\
ROLL.NUMBER.THREE}
\and
\IEEEauthorblockN{Author Four}
\IEEEauthorblockA{\textit{Department} \\
\textit{University}\\
City, Country \\
ROLL.NUMBER.FOUR}
}

\maketitle

\begin{abstract}
Universities must decide how many sections (teaching groups) to open for every course
before an academic year begins. Opening too few sections leaves students unable to
enrol in courses they need; opening too many wastes teaching staff, rooms and timetable
slots. In this paper we formulate and implement an Operations Research decision-support
framework for university curriculum flow and capacity planning. Using the publicly
available StudentDropoutDataset (Zenodo, CC~BY~4.0), comprising 464{,}739 student--course
enrolment records covering 39{,}364 students, 4{,}989 courses and 163 degree programmes
across the academic years 2018, 2021 and 2022, we first project next-year seat demand
per course with a validated student- and course-flow model, and then allocate sections
with two integer optimisation models. Model~1 is an Integer Programming (IP) formulation
that minimises the total number of sections subject to seating all projected demand.
Model~2 is a Binary (0--1) Programming formulation of knapsack type that maximises the
number of students seated subject to a fixed budget of additional sections, an ordering
constraint and an optional per-programme fairness cap. Both models are solved with the
HiGHS solver through the SciPy optimisation interface. At a planning capacity of 50
students per section, the flow model projects 161{,}620 seats of demand for 2023, of
which 68 courses in 28 programmes are bottlenecks and 436 students would not obtain a
seat. Model~1 returns a proven-optimal allocation of 5{,}612 sections, a net change of
$+42$ sections against the 5{,}570 currently modelled (67 courses gain sections, 25
release them). Model~2 reduces unmet demand from 436 to 9 students, a 98.0\,\% reduction,
using a budget of only 50 additional sections; the linear-programming relaxation gives a
budget shadow price of 0.922 students seated per additional section. The framework is
deployed as an interactive Flask web application, verified against a hand calculation and
an Excel Solver workbook on a reduced instance, and covered by 32 automated tests.
\end{abstract}

\begin{IEEEkeywords}
Curriculum planning, capacity planning, integer programming, binary programming,
knapsack problem, operations research, resource allocation, shadow price, HiGHS solver,
educational data.
\end{IEEEkeywords}

\section{Introduction}

University timetabling and capacity planning are resource-allocation problems with a
hard integrality structure. A section is an indivisible teaching unit: a department can
open three sections of a course or four, but not 3.4. Each section consumes a lecturer,
a room and a timetable slot, so sections are scarce and must be justified. Planners must
nevertheless guarantee that students can enrol in the courses their degree requires,
because a student blocked from a required course is delayed by a full academic term.

Capacity decisions are complicated by the fact that demand for a course is not simply
last year's enrolment. Demand has two components: students reaching the course for the
first time through normal progression, and students retaking the course after failing
it. The second component is driven by course-specific pass rates and is systematically
underestimated when planners extrapolate headcounts directly.

Operations Research provides the appropriate machinery \cite{hillier}. Educational resource allocation has a long literature in automated timetabling \cite{burke}, which largely addresses the assignment of events to rooms and slots once section counts are known; the complementary question of how many sections to open is the subject of this paper. When the planning budget is
unconstrained, the question ``what is the smallest number of sections that still seats
every student?'' is an Integer Programming problem. When the budget is fixed in advance,
as it usually is in practice, the question becomes ``given $B$ additional sections,
where do they do the most good?'', which is a Binary (0--1) knapsack problem. The two
formulations answer genuinely different managerial questions and we present both.

The principal contributions of this paper are:
\begin{enumerate}
\item A reproducible preprocessing and validation pipeline for a public, anonymised
      university enrolment dataset of 464{,}739 records, with every published record
      count independently reproduced.
\item A student- and course-flow model that projects next-year seat demand per course,
      separating first-time from repeat demand, validated by a split-half procedure at
      4.2\,\% mean absolute percentage error on student stocks and $R^2 = 0.91$ on
      per-course repeat demand.
\item An Integer Programming formulation (Model~1) that minimises total sections subject
      to full demand coverage, together with its LP relaxation and integrality gap.
\item A Binary (0--1) Programming formulation (Model~2) of knapsack type that maximises
      students seated under a fixed section budget, an ordering constraint on successive
      sections of the same course, and an optional per-programme fairness cap.
\item Duality and sensitivity analysis, including the shadow price of the section budget,
      the marginal value of each successive section, and sensitivity of the solution to
      capacity, budget and intake growth.
\item A three-way verification protocol in which the same reduced instance is solved by
      hand, in Excel Solver and by the deployed application, with agreement required.
\item Deployment as an interactive Flask web application with a dashboard, dataset
      explorer, a page per model and a model-comparison page, supported by 32 automated
      tests.
\end{enumerate}

A deliberate scope restriction applies throughout. Section capacity is a
\emph{planning parameter chosen by the university}, not a variable observed in the
dataset; the published data contains no room or capacity field. Furthermore, this work
does \emph{not} predict which individual students will drop out. The models allocate
teaching capacity; they do not profile students.

\section{Study Data and Planning Parameters}

\subsection{Empirical Database Overview}

The empirical foundation is the StudentDropoutDataset published on Zenodo under the
Creative Commons Attribution 4.0 licence \cite{zenodo}. The archive contains
student--course enrolment records from a Spanish public university for the academic years
2018, 2021 and 2022. Identifiers for students, courses and programmes are published as
irreversible hashes; no personally identifying information is present. Hashes serve only
as join keys, so they constrain neither the arithmetic nor the interpretation of the
models.

\subsection{Preprocessing and Count Validation}

Records are cleaned, typed and linked into a student--course table and a student--year
table. Because the Zenodo record publishes no column codebook, every derived field is
validated against the record counts stated in the dataset description. Table~\ref{tab:valid}
reports this reconciliation; all five published counts are reproduced exactly, which is
the evidence that the column interpretation is correct.

\begin{table}[htbp]
\caption{Dataset validation against published record counts}
\label{tab:valid}
\centering
\begin{tabular}{lrrc}
\toprule
Quantity & Expected & Observed & Match \\
\midrule
Enrolment records & 464{,}739 & 464{,}739 & Yes \\
Students          & 39{,}364  & 39{,}364  & Yes \\
Courses           & 4{,}989   & 4{,}989   & Yes \\
Programmes        & 163       & 163       & Yes \\
Academic years    & 3         & 3         & Yes \\
\bottomrule
\end{tabular}
\end{table}

Table~\ref{tab:eda} summarises the cleaned data by academic year. The pass rate declines
slightly across the observed years, from 73.9\,\% in 2018 to 71.7\,\% in 2022. In 2022,
13{,}446 enrolments (8.4\,\% of all enrolments) were retakes, and 5{,}474 students
(26.8\,\% of the student body) retook at least one course. Repeat enrolment is therefore
a material component of demand and is modelled explicitly.

\begin{table}[htbp]
\caption{Descriptive statistics by academic year}
\label{tab:eda}
\centering
\small
\begin{tabular}{lrrrrr}
\toprule
Year & Records & Students & Courses & Pass rate & Mean grade \\
\midrule
2018 & 152{,}446 & 20{,}233 & 3{,}747 & 73.9\,\% & 6.39 \\
2021 & 153{,}120 & 20{,}264 & 3{,}859 & 72.9\,\% & 6.41 \\
2022 & 159{,}173 & 20{,}427 & 3{,}787 & 71.7\,\% & 6.34 \\
\bottomrule
\end{tabular}
\end{table}

\subsection{Planning Parameters}

The parameters in Table~\ref{tab:params} are assumptions of the planning model. They are
set by the planner and exposed as controls in the web application; none of them is read
from the dataset.

\begin{table}[htbp]
\caption{Planning parameters (model assumptions, not dataset variables)}
\label{tab:params}
\centering
\begin{tabular}{llr}
\toprule
Symbol & Parameter & Value \\
\midrule
$cap$   & Section capacity (students per section) & 50 \\
$g$     & First-time demand growth per year       & 0\,\% \\
$B$     & Additional sections available (Model~2) & 50 \\
$M_p$   & Maximum extra sections per programme    & 5 \\
$\theta$& Bottleneck threshold (demand / capacity)& 1.00 \\
$T$     & Planning horizon for scenarios (years)  & 5 \\
\bottomrule
\end{tabular}
\end{table}

Because the dataset contains no capacity field, the number of sections currently
operating is itself modelled. For each course $c$ the current section count is taken as
\begin{equation}
s_c = \left\lceil \frac{E_c}{cap} \right\rceil ,
\end{equation}
where $E_c$ is observed 2022 enrolment. At $cap = 50$ this yields 5{,}570 sections in
total, which is the baseline against which both models are compared.

\subsection{Demand Projection}

Demand is projected with a stock-and-flow formulation in the system-dynamics tradition \cite{forrester}, in which students accumulate in year-level stocks and move between them at estimated rates.

Let $F_c$ denote projected first-time demand for course $c$, $\rho_c$ the course-specific
retake rate estimated by empirical-Bayes shrinkage toward the university-wide rate, and
$R_c$ the pool of students eligible to retake. Projected demand for the planning year is
\begin{equation}
D_c = F_c (1 + g) + \rho_c R_c .
\label{eq:demand}
\end{equation}

Applying \eqref{eq:demand} with $g = 0$ gives a total projected demand of 161{,}620
seats for 2023, decomposing into 145{,}727 first-time seats and 15{,}893 repeat seats.

The projection is validated by a split-half procedure: flow rates are estimated on one
random half of the student population and used to predict the other half. The mean
absolute percentage error across the student-flow stocks is 4.2\,\%. For per-course
repeat demand, course-specific retake rates achieve $R^2 = 0.909$ with a mean absolute
error of 0.77 students per course, against $R^2 = 0.743$ for a single university-wide
rate. The course-specific estimator carries a small upward bias of $+6.8\,\%$ in total,
which is conservative for a capacity-planning application because it errs toward
providing seats.

\section{Mathematical Formulations}

Let $i \in \{1,\dots,n\}$ index the $n = 3{,}787$ courses offered in the planning year,
with projected demand $D_c$ from \eqref{eq:demand} and section capacity $cap$.

\subsection{Model 1: Minimum Section Allocation (Integer Programming)}

\textbf{Decision variable.} Let $x_c \in \mathbb{Z}_+$ denote the number of sections
opened for course $c$ \cite{nemhauser}. Sections are indivisible, which is what makes this an integer
programme rather than a linear one.

\textbf{Objective.} Minimise the total number of sections opened:
\begin{equation}
\min \; Z_1 = \sum_{c=1}^{n} x_c .
\label{eq:m1obj}
\end{equation}

\textbf{Constraints.}
\begin{align}
cap \cdot x_c &\ge (1+\beta) D_c, & c &= 1,\dots,n, \label{eq:m1cover}\\
x_c &\ge 1, & \forall\, c : D_c > 0, \label{eq:m1min}\\
x_c &\in \mathbb{Z}_+, & c &= 1,\dots,n. \label{eq:m1int}
\end{align}

Constraint \eqref{eq:m1cover} is the service guarantee: the seats provided must cover
projected demand, optionally inflated by a safety buffer $\beta \ge 0$ (set to zero in
the reported results). Constraint \eqref{eq:m1min} ensures that any course with positive
demand actually runs.

A structural property of this formulation is worth stating because it has practical
consequences. Each constraint in \eqref{eq:m1cover} involves only its own variable
$x_c$, so the problem \emph{separates by course} and the optimum has the closed form
\begin{equation}
x_c^{\star} = \max\!\left(1, \left\lceil \frac{(1+\beta) D_c}{cap} \right\rceil \right).
\label{eq:m1closed}
\end{equation}
Model~1 is therefore solvable by hand, which we exploit in
Section~\ref{sec:verification} as an exact check on the solver. Model~2 has no such
property.

\subsection{Model 2: Allocation of Limited Additional Capacity (Binary Programming)}

In practice the number of extra sections that can be funded is fixed in advance. Let
$s_c$ be the sections currently operating and define the unmet demand of course $c$ as
\begin{equation}
U_c = \max\!\left(0,\; D_c - cap \cdot s_c\right).
\end{equation}

\textbf{Decision variable.} For course $c$, let $y_{ck} \in \{0,1\}$ equal $1$ if the
$k$-th additional section is opened for that course and $0$ otherwise, for
$k = 1,\dots,K_c$ with $K_c = \lceil U_c / cap \rceil$. Each candidate section is an item
and the budget is the knapsack \cite{martello}.

\textbf{Coefficients.} The $k$-th additional section of course $c$ seats
\begin{equation}
a_{ck} = \min\!\big(cap,\; U_c - (k-1)\,cap\big)
\end{equation}
students, which is a diminishing quantity: the last section opened for a course is
usually only partly filled. Let $w_c \ge 1$ be a course priority weight
(Section~\ref{sec:ahp}); $w_c = 1$ for all $c$ in the reported results.

\textbf{Objective.} Maximise the weighted number of students seated:
\begin{equation}
\max \; Z_2 = \sum_{c=1}^{n} \sum_{k=1}^{K_c} w_c\, a_{ck}\, y_{ck}.
\label{eq:m2obj}
\end{equation}

\textbf{Constraints.}
\begin{align}
\sum_{c=1}^{n} \sum_{k=1}^{K_c} y_{ck} &\le B, \label{eq:m2budget}\\
y_{ck} &\le y_{c,k-1}, \quad k = 2,\dots,K_c, \label{eq:m2order}\\
\sum_{c \in p} \sum_{k=1}^{K_c} y_{ck} &\le M_p, \quad \forall\, \text{programmes } p,
\label{eq:m2fair}\\
y_{ck} &\in \{0,1\}. \label{eq:m2bin}
\end{align}

Constraint \eqref{eq:m2budget} is the section budget. Constraint \eqref{eq:m2order}
forces sections of a course to be opened in order, which prevents the solver from
claiming the high coverage of a first section while nominally opening a later one.
Constraint \eqref{eq:m2fair} is an optional fairness cap limiting how many extra sections
any one programme may absorb. Maximising coverage in \eqref{eq:m2obj} is equivalent to
minimising total residual unmet demand $\sum_c U_c - \sum_{c,k} a_{ck} y_{ck}$.

Unlike Model~1, Model~2 does \emph{not} separate by course: all courses compete for the
shared budget in \eqref{eq:m2budget}, and the fairness cap couples courses within a
programme. This coupling is what genuinely requires an optimisation solver.

\subsection{LP Relaxation and Sensitivity Analysis}

Relaxing $y_{ck} \in \{0,1\}$ to $0 \le y_{ck} \le 1$ yields a linear programme whose
dual provides managerial information \cite{dantzig}. The shadow price of \eqref{eq:m2budget} is the
number of additional students seated by one more section of budget. Because each
$y_{ck} \le 1$, the objective is piecewise linear in $B$ with a breakpoint at every whole
section, so the ranging table reports the marginal value of each successive section.
The LP relaxation of Model~1 is likewise solved to obtain a lower bound on
\eqref{eq:m1obj} and hence the integrality gap.

\subsection{Course Priority Weights via AHP}
\label{sec:ahp}

When the planner wishes to express a priority other than raw headcount, the weights
$w_c$ in \eqref{eq:m2obj} may be derived by the Analytic Hierarchy Process \cite{saaty} over three
criteria: the repeat-student share of a course's demand, whether the course is a
first-year course, and one minus its pass rate. For a representative pairwise comparison
matrix the resulting criteria weights are $0.648$, $0.230$ and $0.122$ respectively,
with a consistency ratio of $CR = 0.003$, well inside the conventional $CR \le 0.10$
acceptance threshold. The course weight is then
$w_c = 1 + \sum_k W_k \cdot \mathrm{score}_k(c)$.

\section{System Architecture and Implementation}

The system is implemented in Python 3.13 with a modular architecture that separates the
analytical core from the interface:

\begin{itemize}
\item \texttt{preprocessing.py} --- loads the published archive, cleans and links the
      student--course tables, and writes the count-validation log.
\item \texttt{system\_dynamics.py} --- student- and course-flow models, empirical-Bayes
      rate shrinkage, and the split-half validation routines.
\item \texttt{capacity.py} --- modelled current sections, bottleneck classification and
      bottleneck summaries.
\item \texttt{optimization.py} --- constructs the constraint matrices for both models and
      calls \texttt{scipy.optimize.milp} \cite{scipy}, which uses the HiGHS branch-and-bound and
      dual-simplex solvers \cite{highs}; extracts primal solutions, slacks and duals.
\item \texttt{mcdm.py} --- AHP pairwise comparison, principal eigenvector weights, and
      the consistency ratio.
\item \texttt{scenarios.py} --- couples the flow model and the optimisation models in a
      multi-year loop, with shortfall returning as next-year backlog, and performs the
      sensitivity sweeps.
\item \texttt{run\_pipeline.py} --- batch execution writing all tables, figures and the
      auto-generated results report.
\item \texttt{flask\_app.py} --- Flask web application with a dashboard, a dataset
      explorer providing statistics on demand, a page for each optimisation model, and a
      model-comparison page; charts are rendered with Chart.js.
\item Test suite --- 32 automated tests covering unit behaviour of the formulations and
      end-to-end assertions against the real dataset.
\end{itemize}

Heavy computation (preprocessing, flow calibration, demand projection) is performed once
by the pipeline and cached; the web application re-solves both optimisation models live
for any parameter setting, which completes in well under one second for the full
3{,}787-course instance.

\section{Results and Discussion}

\subsection{Baseline and Bottleneck Diagnosis}

At $cap = 50$ the modelled current provision is 5{,}570 sections, giving 278{,}500 seats
against 161{,}620 seats of projected demand. Aggregate capacity is therefore ample; the
problem is distributional. Of 3{,}787 courses, 68 are bottlenecks
($D_c \ge \theta \cdot cap \cdot s_c$) and a further 267 are at risk at the 90\,\%
utilisation threshold. The bottlenecks span 28 degree programmes and leave 436 students
without a seat. Repeat enrolment accounts for 16.4\,\% of demand in bottleneck courses,
against 9.8\,\% across all courses, confirming that retakes concentrate in exactly the
courses that are already full.

\subsection{Model 1: Minimum Section Allocation}

Model~1 returns a proven-optimal allocation of 5{,}612 sections (HiGHS status: optimal;
all variables integral; demand fully met). Against the 5{,}570 modelled current sections
this is a net change of only $+42$: 67 courses require additional sections and 25 courses
can release a section without breaching \eqref{eq:m1cover}. The managerial reading is
that the institution is not short of teaching capacity in aggregate but is misallocating
it, and that full coverage is achievable at close to budget neutrality.

The LP relaxation of \eqref{eq:m1obj}--\eqref{eq:m1int} attains 5{,}062.3 fractional
sections, so the integrality gap is 549.7 sections. The gap is large because
\eqref{eq:m1min} forces a whole section onto every small course; this also explains the
57.6\,\% mean seat utilisation of the integer solution, which is a consequence of
indivisibility rather than of over-provision.

\subsection{Model 2: Limited Additional Capacity}

With a budget of $B = 50$ additional sections and uniform weights, Model~2 seats 427.7
of the 436.5 unmet students, reducing unmet demand to 8.8 students, a 98.0\,\%
reduction. The number of courses still short falls from 67 to 17. All 50 sections are
used, each going to a distinct course. Clearing the shortfall entirely requires 67
sections. Table~\ref{tab:compare} compares the baseline against both models.

\begin{table}[htbp]
\caption{System performance comparison: baseline vs.\ Model 1 vs.\ Model 2}
\label{tab:compare}
\centering
\small
\begin{tabular}{lrrr}
\toprule
Evaluation metric & Baseline & Model~1 & Model~2 \\
\midrule
Sections operated        & 5{,}570 & 5{,}612 & 5{,}620 \\
Change vs.\ baseline     & ---     & $+42$   & $+50$ \\
Students without a seat  & 436.5   & 0.0     & 8.8 \\
Reduction in unmet demand& ---     & 100.0\,\% & 98.0\,\% \\
Courses still short      & 67      & 0       & 17 \\
Seats provided           & 278{,}500 & 280{,}600 & 281{,}000 \\
Mean seat utilisation    & 58.0\,\% & 57.6\,\% & 57.5\,\% \\
Solver status            & ---     & Optimal & Optimal \\
\bottomrule
\end{tabular}
\end{table}

The two models answer different questions and neither dominates. Model~1 describes the
ideal allocation when the budget is open; Model~2 describes the best achievable outcome
when it is not. Model~2 attains 98\,\% of the available benefit with a budget 25\,\%
smaller than the 67 sections that full coverage requires, which is the practically
relevant finding.

\subsection{Duality, Binding Constraints and Shadow Prices}

The LP relaxation of Model~2 yields the following dual information at $B = 50$:

\begin{itemize}
\item \textbf{Section budget.} The budget constraint \eqref{eq:m2budget} is strictly
      binding (50 of 50 sections used). Its shadow price is \textbf{0.922 students per
      section}: one additional section of budget seats approximately 0.92 further
      students at this point in the ranging table.
\item \textbf{Marginal value of successive sections.} The 50th section is worth 0.938
      students and the 51st is worth 0.922, confirming diminishing returns. Early
      sections are far more valuable: the first five sections seat 30.3 students each on
      average, falling to 15.1 at $B=10$, 8.4 at $B=25$ and 3.0 at $B=50$.
\item \textbf{Saturation.} Beyond $B = 67$ the shadow price falls to zero. No further
      section can seat anyone because all unmet demand has been absorbed; additional
      budget is pure slack.
\item \textbf{Fairness cap.} Imposing $M_p = 5$ extra sections per programme raises
      residual unmet demand from 8.8 to 19.6 students. The cost of the equity
      constraint is therefore 10.8 students, which quantifies the efficiency--equity
      trade-off explicitly rather than leaving it to judgement.
\end{itemize}

\subsection{Sensitivity Analysis}

Table~\ref{tab:budget} reports sensitivity to the section budget and
Table~\ref{tab:capacity} to the capacity parameter.

\begin{table}[htbp]
\caption{Sensitivity to the additional-section budget (Model 2)}
\label{tab:budget}
\centering
\small
\begin{tabular}{rrrr}
\toprule
Budget $B$ & Sections used & Unmet after & Marginal coverage \\
\midrule
0   & 0  & 436.5 & --- \\
5   & 5  & 285.0 & 30.29 \\
10  & 10 & 209.4 & 15.11 \\
25  & 25 & 83.9  & 8.37 \\
50  & 50 & 8.8   & 3.00 \\
75  & 67 & 0.0   & 0.35 \\
100 & 67 & 0.0   & 0.00 \\
\bottomrule
\end{tabular}
\end{table}

\begin{table}[htbp]
\caption{Sensitivity to the section-capacity parameter}
\label{tab:capacity}
\centering
\small
\begin{tabular}{rrrrr}
\toprule
$cap$ & Bottlenecks & At risk & Model~1 sections & Utilisation \\
\midrule
30  & 133 & 472 & 7{,}533 & 71.5\,\% \\
40  & 86  & 345 & 6{,}278 & 64.4\,\% \\
50  & 68  & 267 & 5{,}612 & 57.6\,\% \\
60  & 48  & 206 & 5{,}200 & 51.8\,\% \\
75  & 38  & 202 & 4{,}752 & 45.3\,\% \\
100 & 31  & 128 & 4{,}337 & 37.3\,\% \\
\bottomrule
\end{tabular}
\end{table}

Larger sections reduce both the number of bottlenecks and the number of sections
required, but at a monotonic cost in seat utilisation and, in practice, in teaching
quality. The capacity parameter is therefore a genuine policy lever rather than a
technical constant, which is precisely why it is exposed in the interface.

A caveat applies to Table~\ref{tab:capacity}. Status-quo unmet demand is \emph{not}
monotone in $cap$, because the modelled baseline $s_c = \lceil E_c / cap \rceil$ is
itself recomputed at each capacity. Rounding effects therefore move the baseline as the
parameter moves. The Model~1 column, which does not depend on the baseline, is monotone
as expected.

\subsection{Multi-Year Scenario Analysis}

Coupling the flow model with each allocation policy over a five-year horizon, with any
shortfall returning as next-year backlog, gives Table~\ref{tab:scen}.

\begin{table}[htbp]
\caption{Scenario comparison, 2023--2027}
\label{tab:scen}
\centering
\small
\begin{tabular}{llrr}
\toprule
Scenario & Policy & Year-1 unserved & Cumulative \\
\midrule
S0 Status quo              & none     & 436.5 & 3{,}389.5 \\
S1 Limited extra           & Model~2, $B{=}50$  & 8.8 & 94.5 \\
S2 Double budget           & Model~2, $B{=}100$ & 0.0 & 0.0 \\
S3 Full coverage           & Model~1  & 0.0 & 0.0 \\
S4 Model~2 $+$ pass $+5$pp & Model~2  & 8.8 & 12.8 \\
S5 Model~2, intake $+5\%$/yr & Model~2 & 354.9 & 29{,}344.6 \\
\bottomrule
\end{tabular}
\end{table}

Scenario S5 is the substantive result. A sustained 5\,\% annual growth in first-time
intake drives cumulative unserved demand to 29{,}344 students over five years despite
the same 50-section intervention that is almost sufficient under flat intake. No
realistic section budget absorbs this: the binding issue becomes intake policy rather
than timetable capacity. Conversely, S4 shows that a five-percentage-point improvement
in pass rates reduces cumulative unserved demand from 94.5 to 12.8, because fewer
retakes are generated, and simultaneously reduces the terminal backlog from 46{,}348 to
37{,}454. Academic support and capacity planning are complementary interventions.

\subsection{Three-Way Verification}
\label{sec:verification}

To validate the implementation independently of the solver, a reduced instance of the
eight largest bottleneck courses was solved by three methods: by hand, in Microsoft Excel
Solver (Simplex LP with integer and binary constraints), and by the deployed application.
Table~\ref{tab:verify} reports the outcome. Agreement is exact.

\begin{table}[htbp]
\caption{Three-way verification on the eight-course instance}
\label{tab:verify}
\centering
\small
\begin{tabular}{llr}
\toprule
Model & Method & Result \\
\midrule
Model~1 & Hand, eq.~\eqref{eq:m1closed} & 48 sections \\
Model~1 & Excel Solver (integer)         & 48 sections \\
Model~1 & Application (HiGHS)            & 48 sections \\
\midrule
Model~2 & Hand, $B{=}3$, no cap          & 112.22 students \\
Model~2 & Hand, $B{=}3$, $M_p{=}1$       & 88.62 students \\
Model~2 & Application (HiGHS), $M_p{=}1$ & 88.62 students \\
\bottomrule
\end{tabular}
\end{table}

The Model~2 rows demonstrate the non-separability claimed in Section~III. Without the
fairness cap the optimum is obtained by selecting the three courses with the largest
unmet demand. With $M_p = 1$ that greedy selection becomes infeasible, the optimal
selection changes, and the objective falls by 23.60 students. A greedy heuristic would
return the wrong answer; the binary programme does not.

\subsection{Effect of Priority Weights}

An honest negative result should be reported. Substituting repeat-weighted priorities for
uniform weights in \eqref{eq:m2obj} leaves the solution unchanged at $B = 50$
(8.8 students unmet in both cases), and AHP-derived weights likewise. The coverage
coefficients $a_{ck}$ dominate the weight multipliers at this budget, so the weighting
scheme is expressive but not, on this dataset, decisive. We report the AHP machinery and
its consistency check as a methodological component rather than as a driver of the
result.

\section{Limitations}

\begin{enumerate}
\item \textbf{Capacity is assumed homogeneous.} Every section is assumed to seat $cap$
      students. Real rooms differ in size, and laboratory or studio courses have
      different ratios. The model treats capacity as a single policy parameter.
\item \textbf{No timetable feasibility.} The models determine how many sections to open,
      not when they meet or who teaches them. A section count that is feasible in
      aggregate may be infeasible once room and staff availability are imposed.
\item \textbf{Deterministic demand.} Projected demand is a point estimate. Stochastic
      variation in enrolment is not propagated into the allocation, although the safety
      buffer $\beta$ provides a deterministic proxy.
\item \textbf{Non-consecutive observation years.} The archive provides 2018, 2021 and
      2022, so only the 2021--2022 pair is consecutive and available for flow
      calibration; the 2018 cohort informs descriptive statistics only.
\item \textbf{Undocumented schema.} The Zenodo record publishes no column codebook.
      Column interpretation rests on the exact reproduction of all published record
      counts (Table~\ref{tab:valid}) and on internal consistency checks.
\end{enumerate}

\section{Conclusion and Future Work}

This paper presented an Operations Research decision-support system for university
curriculum flow and capacity planning. Using 464{,}739 enrolment records from a public,
anonymised dataset, a validated flow model projects 161{,}620 seats of demand for the
planning year, of which 68 courses are bottlenecks leaving 436 students without a seat.

An Integer Programming formulation establishes that full coverage requires 5{,}612
sections, a net increase of only 42 over the 5{,}570 currently modelled, because 67
courses gain sections while 25 release them. A Binary (0--1) Programming formulation
establishes that a fixed budget of 50 additional sections removes 98.0\,\% of the
shortfall, reducing unmet demand from 436 to 9 students, with a budget shadow price of
0.922 students per section and complete saturation at 67 sections. Imposing a
per-programme fairness cap costs 10.8 students, quantifying the equity--efficiency
trade-off. Both models solve to proven optimality and agree exactly with hand and Excel
Solver computations on a reduced instance.

The scenario analysis qualifies the recommendation. Under flat intake the proposed
intervention is close to sufficient, but a sustained 5\,\% annual growth in first-time
intake generates 29{,}344 cumulative unserved students over five years, which no
plausible section budget absorbs. Capacity optimisation is necessary but not sufficient;
it must be paired with intake planning and with measures that raise pass rates and so
reduce retake demand.

Future work will extend the formulation in three directions: a two-stage stochastic
programme in which the second-stage recourse reallocates sections after observed
enrolment; a multi-period model in which section decisions in one year constrain the
next through staff contracts; and integration of room-level capacity and timetable
feasibility so that the allocation produced is directly implementable.

\section*{Reproducibility}

The dataset is publicly available under CC~BY~4.0 \cite{zenodo}. All results in this
paper are regenerated end to end by a single command and are accompanied by 32 automated
tests. Source code and generated outputs are available at
\url{https://github.com/gopinathgatta/OR_project}.

\begin{thebibliography}{00}

\bibitem{hillier}
F. S. Hillier and G. J. Lieberman, \emph{Introduction to Operations Research}, 11th ed.
New York, NY, USA: McGraw-Hill Education, 2021.

\bibitem{dantzig}
G. B. Dantzig, \emph{Linear Programming and Extensions}. Princeton, NJ, USA: Princeton
University Press, 1963.

\bibitem{nemhauser}
G. L. Nemhauser and L. A. Wolsey, \emph{Integer and Combinatorial Optimization}.
New York, NY, USA: Wiley-Interscience, 1988.

\bibitem{martello}
S. Martello and P. Toth, \emph{Knapsack Problems: Algorithms and Computer
Implementations}. Chichester, U.K.: John Wiley \& Sons, 1990.

\bibitem{saaty}
T. L. Saaty, ``How to make a decision: The Analytic Hierarchy Process,''
\emph{European Journal of Operational Research}, vol. 48, no. 1, pp. 9--26, 1990.

\bibitem{highs}
Q. Huangfu and J. A. J. Hall, ``Parallelizing the dual revised simplex method,''
\emph{Mathematical Programming Computation}, vol. 10, no. 1, pp. 119--142, 2018.

\bibitem{scipy}
P. Virtanen \emph{et al.}, ``SciPy 1.0: Fundamental algorithms for scientific computing
in Python,'' \emph{Nature Methods}, vol. 17, pp. 261--272, 2020.

\bibitem{burke}
E. K. Burke and S. Petrovic, ``Recent research directions in automated timetabling,''
\emph{European Journal of Operational Research}, vol. 140, no. 2, pp. 266--280, 2002.

\bibitem{forrester}
J. W. Forrester, \emph{Industrial Dynamics}. Cambridge, MA, USA: MIT Press, 1961.

\bibitem{zenodo}
StudentDropoutDataset, Zenodo, 2025. [Online]. Available:
\url{https://doi.org/10.5281/zenodo.17239943}

\end{thebibliography}

\end{document}
